\chapter{Introduction}

\section{Background and Motivation}

AWS and Microsoft Azure now each offer well over two hundred distinct services \cite{aws2025catalog, azure2025catalog}, a figure that has roughly doubled over the past decade \cite{flexera2024cloud}. This expansion gives architects unprecedented flexibility, but it also compounds decision complexity: selecting and integrating the right services for a given workload requires weighing competing factors---cost, scalability, security, and performance---across offerings that frequently overlap in purpose but differ substantially in API design and operational characteristics.

This thesis proposes an agentic AI framework to assist architects in navigating this architectural complexity. The framework employs an Evaluator-Optimizer workflow \cite{anthropic2024agents} that iteratively generates, validates, and refines cloud architectural recommendations. Validation is grounded in authoritative cloud provider documentation through Retrieval Augmented Generation (RAG) \cite{lewis2020rag, asai2023selfrag}, ensuring factual and reliable outputs. The system aims to deliver efficient, secure, and cost-optimized solutions capable of cross-provider comparative analysis. %\cite{futransolutions2025multi-cloud}.

\section{Problem Statement}

The transition from traditional on-premises infrastructure to cloud platforms has resulted in an exponential increase in cloud services. Providers like AWS and Azure often offer overlapping services solving similar problems, complicating architectural decisions. %\cite{simplilearn2025.aws-vs-azure-challenges}
Architects typically rely on domain expertise and experience to evaluate tradeoffs involving costs, scalability, security, and performance. Current AI-based tools and assistants, including first-party vendor AI products like Amazon Q \cite{aws2023q} and Microsoft Copilot \cite{microsoft2024copilot}, suffer from several limitations:

\begin{itemize}
    \item Existing chatbots often provide partial guidance focused on isolated components rather than end-to-end architecture coverage.
    \item Generated outputs frequently lack grounded factual validation against authoritative cloud documentation.
    \item Current solutions generally do not support comparative evaluation across multiple cloud providers, often leading to vendor-biased recommendations.
\end{itemize}

These gaps motivate an AI-driven framework that generates validated, end-to-end architectural recommendations and supports side-by-side multi-cloud comparison to reduce vendor bias in decision-making. %\cite{agiliway2025aws-vs-azure, futransolutions2025multi-cloud}.

\section{Objectives}

The primary objectives of this research project are divided into system architecture, practical implementation, and formal evaluation:

\subsubsection*{System Architecture and AI Workflow}
\begin{itemize}
    \item To design and develop autonomous AI agent teams representing leading cloud providers (AWS and Azure) to generate architecture proposals.
    \item To design an Evaluator-Optimizer workflow using Retrieval-Augmented Generation (RAG) and document reranking for the factual validation of recommendations.
    \item To enable a multi-turn agent debate mechanism that facilitates side-by-side comparative evaluations between AWS and Azure architectures.
\end{itemize}

\subsubsection*{Implementation and Deployment}
\begin{itemize}
    \item To provide Infrastructure-as-Code (IaC) generation capabilities (Terraform, CloudFormation, Bicep) to translate theoretical designs into directly deployable artifacts.
    \item To build and deploy a functional, end-to-end prototype system accessible via user interface.
\end{itemize}

\subsubsection*{Validation and Evaluation}
\begin{itemize}
    \item To validate the semantic quality and accuracy of the generated solutions utilizing established lexical similarity metrics (e.g., METEOR, BERTScore).
    \item To evaluate framework quality using an LLM-as-a-judge protocol deployed across stringent technical dimensions.
\end{itemize}


\section{Scope and Limitations}

\subsection{Scope}

The research focuses on developing and evaluating a functional prototype of the agentic AI framework, limited to the following boundaries:

\begin{itemize}
    \item \textbf{Cloud Providers:} Initially limited to Amazon Web Services (AWS) and Microsoft Azure, the two dominant public cloud vendors.
    \item \textbf{Knowledge Source:} The Retrieval Augmented Generation knowledge base is constructed exclusively from publicly available, official provider documentation such as whitepapers, service guides, and recommended best practices.
    \item \textbf{Output Format:} The system outputs comprehensive textual cloud-architecture recommendations with detailed justifications, as well as generates deployable code (Infrastructure-as-Code templates such as Terraform or CloudFormation).
\end{itemize}

\subsection{Limitations}

The implemented prototype currently has the following operational constraints:

\begin{itemize}
    \item \textbf{Absence of Real-Time Data:} The framework does not incorporate real-time information such as pricing updates, service availability in specific regions, or customer-specific account configurations.
    \item \textbf{Knowledge Base Boundaries:} Although extensive, the knowledge base may lack coverage of some niche or recently introduced cloud services due to document availability and constraints associated with the open source vector database.
    \item \textbf{Non-Substitutive Advisory Role:} The system is intended as a decision support tool complementing human expertise rather than replacing human operators. All architectural recommendations warrant expert review prior to deployment.
\end{itemize}

\section{Feasibility Study}

\subsection{Technical Feasibility}

Mature underlying technologies such as Large Language Models (LLMs), vector similarity searches, and agentic AI frameworks (e.g., LangChain \cite{langchain2024docs}, LangGraph \cite{langgraph2024docs}, AutoGen \cite{wu2023autogen}) make this framework technically realisable. Cloud platforms (AWS and Azure) provide rich documentation and resources essential for building the knowledge base.

\subsection{Data Availability}

Cloud vendor documentation required for the RAG knowledge base is openly accessible. Such resources include service whitepapers, architecture best practices, and official guides, forming a sufficient foundation for information retrieval and validation. %\cite{aws2025well-architected, azure2025-architecture}.

\subsection{Operational Feasibility}

The prototype targets deployment on affordable local computation platforms using GPU-accelerated frameworks (Apple Metal Performance Shaders \cite{apple2024mps}) and local embeddings (Ollama \cite{ollama2024web}), with the potential to integrate cloud based services. The user interfaces will be minimal yet practical comprising command-line and web portals, to streamline accessibility, including for evaluation purposes.

\subsection{Expertise and Challenges}

The primary engineering challenge is ensuring factual accuracy in AI-generated recommendations---a problem addressed through the two-stage RAG pipeline and the iterative Evaluator-Optimizer loop described in Chapters~3 and~4.

\section{Thesis Organization}

The dissertation is organized into the following chapters:

\begin{itemize}
    \item \textbf{Chapter 2} surveys foundational concepts covering evolution of agentic multi-agent systems, Retrieval-Augmented Generation (RAG) with cross-encoder reranking, multi-agent debate paradigms, and the LLM-as-a-judge evaluation methodology. The chapter concludes by identifying the research gaps this work addresses.
    
    \item \textbf{Chapter 3} details the system design and core methodology. It transitions from theory to architecture, explicitly mapping out the conceptual agentic AI frameworks, the theoretical orchestration of the Evaluator-Optimizer workflow, and the overarching algorithmic approach.
    
    \item \textbf{Chapter 4} focuses on the practical implementation and engineering of the system. It breaks down the applied components of the framework, detailing the backend data pipeline (web scraping and Chroma vector database construction), the programmatic realization of the multi-agent workflow (LangGraph), and the frontend user interface deployment (Next.js).
    
    \item \textbf{Chapter 5} presents the empirical evaluations and results. It features an ablation study comparing three framework variants (a baseline, an agentic draft, and the full Evaluator-Optimizer loop). Additionally, it conducts a cross-model performance comparison between leading models (GPT-5.4, Claude, and Gemini~3.1 Pro) using the LLM-as-a-judge protocol, concluding with an analysis of methodological limitations.
    
    \item \textbf{Chapter 6} concludes the thesis by summarizing key findings and contributions. It acknowledges the structural limitations of the current framework and outlines strategic future directions, such as the integration of a Google Cloud Platform (GCP) context and the enabling of live, automated Infrastructure-as-Code (IaC) push deployments.
\end{itemize}